\section*{الفصل \ref{ch:g12:mountain-building} --- بناء الجبال}

\begin{solution}{exo:g12:mountain-building:1}
(أ) \omterm{def:g12:mountain-building:thrust}{الصدع الزاحف}؛ (ب) \omterm{def:g12:mountain-building:ophiolite}{الأوفيوليت}؛ (ج) \omterm{def:g12:mountain-building:fold}{الطية المقعرة}؛ (د) \omterm{def:g12:mountain-building:suture}{الدرز}؛ (ه) \omterm{def:g12:mountain-building:fold}{الطية المحدبة}؛
(و) \omterm{def:g12:mountain-building:nappe}{الغطاء التكتوني}.
\end{solution}

\begin{solution}{exo:g12:mountain-building:2}
(1) يميل الجناحان بعيدًا عن \omterm{def:g12:mountain-building:fold}{محور الطية}: فهي \omterm{def:g12:mountain-building:fold}{طية محدبة}. (2) أقدم
الطبقات. (3) لا: فالتعرية تتبع الطبقات الصلبة والطرية، لا شكل
\omterm{def:g12:mountain-building:fold}{الطية}؛ وقد يتآكل لبّ \omterm{def:g12:mountain-building:fold}{الطية المحدبة} حتى يصير واديًا.
\end{solution}

\begin{solution}{exo:g12:mountain-building:3}
(1) $e = 1 - \cos\ang{30} = 1 - 0.866 = 0.134$، أي نحو \qty{13}{\percent}.
(2) $L_0 = 6.0/\cos\ang{30} \approx \qty{6.9}{km}$.
\end{solution}

\begin{solution}{exo:g12:mountain-building:4}
(1) \omterm{def:g11:ocean-ridges:normal-fault}{صدع عادي}؛ وقد تمططت القشرة. (2) نعم: إنه \omterm{def:g12:mountain-building:thrust}{صدع معكوس}
يميل أقل من \ang{45}.
\end{solution}

\begin{solution}{exo:g12:mountain-building:5}
(1) الرسوبيات البحرية، فالبازلت الوسائدي، \omterm{def:g11:igneous-rocks:intrusion}{فالقواطع} المتراصة، \omterm{def:g11:ocean-ridges:gabbro}{فالغابرو}، \omterm{def:g11:ocean-ridges:peridotite}{فالبيريدوتيت}.
(2) لا تتخذ اللابة هيئة الوسائد إلا إذا بردت في الماء. (3) \omterm{def:g12:mountain-building:ophiolite}{الأوفيوليت}
قطعة من المحيط الذي كان بين قارتين؛ فحين تصادمتا،
دُفعت شرائح من قاعه على حافتيهما بدل أن تغوص، فهي إذن
توجد حيث أُغلق المحيط، على امتداد \omterm{def:g12:mountain-building:suture}{الدرز}.
\end{solution}

\begin{solution}{exo:g12:mountain-building:6}
(1) $L_0 = 2.1 + 1.8 + 2.4 + 1.6 = \qty{7.9}{km}$؛
$e = (7.9 - 6.0)/7.9 \approx 0.24$، أي \qty{24}{\percent}. (2) $\cos\theta =
6.0/7.9 \approx 0.759$، و$\theta \approx \ang{41}$.
\end{solution}

\begin{solution}{exo:g12:mountain-building:7}
(1) $45\cos\ang{10} \approx \qty{44.3}{km}$. (2) $45\sin\ang{10} \approx
\qty{7.8}{km}$. (3) مثل هذه \omterm{def:g12:mountain-building:fold}{الطيات} تقصّر بمقدار $1 - \cos\ang{12} \approx 0.0219$:
فلا بد أن عرض الحزام كان $44.3/0.0219 \approx \qty{2000}{km}$ قبل
الطي.
\end{solution}

\begin{solution}{exo:g12:mountain-building:8}
(1) مبدأ التراكب: فلو قُرئت الصخور دون الصدع، لعُدّ الصخر الأعلى
هو الأحدث. (2) $500 - 75 = 425$ مليون سنة. (3) $\qty{40000}{m} /
\qty{800}{m} = 50$ ضعفًا.
\end{solution}

\begin{solution}{exo:g12:mountain-building:9}
(1) $(30.23 - 19.38) \times 111.2 = 10.85 \times 111.2 \approx
\qty{1207}{km}$. (2) $1207/14 \approx \qty{86}{km}$ في كل مليون سنة،
أي \qty{86}{mm} أو نحو \qty{8.6}{cm} في السنة. (3) \qty{9}{m} في القرن تساوي
\qty{9}{cm} في السنة: فالقيمتان متفقتان في حدود نحو \qty{5}{\percent}.
\end{solution}

\begin{solution}{exo:g12:mountain-building:10}
(1) $H = 35/(1 - 0.40) \approx \qty{58}{km}$. (2) $\Delta H \approx
\qty{23.3}{km}$؛ فهي ترتفع بمقدار $23.3 \times 550/3380 \approx
\qty{3.8}{km}$، وعمق \omterm{def:g12:continental-crust:root}{جذرها} الإضافي $23.3 - 3.8 \approx \qty{19.5}{km}$.
\end{solution}

\begin{solution}{exo:g12:mountain-building:11}
(1) الرسم يضرب $\tan\theta$ في 3: $\tan\theta = \tan\ang{40}/3
\approx 0.280$، و$\theta \approx \ang{15.6}$. (2) الحقيقي: $1 - \cos\ang{15.6}
\approx 0.037$، أي \qty{3.7}{\percent}؛ والمقروء على الرسم: $1 - \cos\ang{40}
\approx 0.234$، أي \qty{23}{\percent}. (3) المقياس الرأسي المُمطَّط يزيد
الميول حدة ويطيل الطبقات، فيجعل \omterm{def:g12:mountain-building:shortening}{التقصير} يبدو أكبر من حقيقته
بعدة أضعاف.
\end{solution}

\begin{solution}{exo:g12:mountain-building:12}
(1) $L_0 = 10/(1 - 0.317) \approx \qty{14.6}{km}$. (2) $1 - \cos\ang{57.5}
\approx 1 - 0.537 = 0.463$، أي \qty{46}{\percent}. (3) في \omterm{def:g12:mountain-building:fold}{الطية} المستديرة
لا يميل بزاوية \ang{57.5} إلا جزء قصير من الطبقة؛ وقرب المفاصل تكون الطبقة
مستوية تقريبًا، والأجزاء المستوية لا تُقصَّر. أما في \omterm{def:g12:mountain-building:fold}{الطية} ذات الجناحين المستقيمين
فالطبقة كلها تميل بزاوية \ang{57.5}.
\end{solution}

\begin{solution}{exo:g12:mountain-building:13}
(1) $\sqrt{16.6^2 + 41.8^2} \approx \qty{45.0}{mm/yr}$. (2) $41.8 - 23.5 =
\qty{18.3}{mm}$ في السنة، أي نحو \qty{18}{km} من \omterm{def:g12:mountain-building:shortening}{التقصير} الشمالي--الجنوبي
في كل مليون سنة شمال لاسا. (3) لاسا نفسها تحركت شمالًا،
وزلازل الحزام تمتد بعيدًا إلى شمال الهيمالايا عبر آسيا: فالقشرة
هناك لا تزال تنكسر.
\end{solution}

\begin{solution}{exo:g12:mountain-building:14}
(1) ميله قريب من الصفر، \omterm{def:g12:mountain-building:shortening}{فالتقصير} يساوي الإزاحة،
\qty{45}{km}. (2) $45/(1 - \cos\ang{12}) = 45/0.0219 \approx \qty{2060}{km}$.
(3) التقاربات الكبيرة تمتصها أساسًا \omterm{def:g12:mountain-building:thrust}{الصدوع الزاحفة} \omterm{def:g12:mountain-building:nappe}{والأغطية التكتونية} \omterm{def:g12:mountain-building:fold}{والطيات}
الشديدة؛ أما \omterm{def:g12:mountain-building:fold}{الطيات} اللطيفة فلا تمتص إلا القليل جدًا.
\end{solution}

\begin{solution}{exo:g12:mountain-building:15}
(1) $\qty{8848.86}{m} = \qty{8848860}{mm}$ في $\num{5e7}$ سنة: أي نحو
\qty{0.18}{mm} في السنة. (2) $40/0.18 \approx 220$: فالارتفاع أبطأ من حركة
الهند بأكثر من مئتي مرة. (3) ترسب الحجر الجيري
تحت مستوى سطح البحر، وقد أزالت التعرية الصخور التي كانت فوقه: فقد
رُفع بأكثر من ارتفاعه الحالي.
\end{solution}

\begin{solution}{pb:g12:mountain-building:1}
\textbf{1.} $2\pi \times 6371/360 \approx \qty{111.2}{km}$.

\textbf{2.} بين 50 مليون سنة مضت (\ang{0.68} جنوبًا) و45 (\ang{3.30} شمالًا):
$50 - 5 \times 0.68/3.98 \approx 49$ مليون سنة مضت.

\textbf{3.} $(19.38 - 0.68) \times 111.2 = 18.70 \times 111.2 \approx
\qty{2079}{km}$.

\textbf{4.} $2079/16 \approx \qty{130}{km}$ في كل مليون سنة: أي \qty{130}{mm}
في السنة، نحو \qty{13}{cm}.

\textbf{5.} من 66 إلى 60: $7.28 \times 111.2/6 \approx 135$؛ ومن 60 إلى 55: $4.86
\times 111.2/5 \approx 108$؛ ومن 55 إلى 50: $6.56 \times 111.2/5 \approx 146$؛ ومن 50
إلى 45: $3.98 \times 111.2/5 \approx 89$؛ ومن 45 إلى 40: $2.29 \times 111.2/5
\approx 51$ (\unit{mm} في السنة). وكانت أسرع من 55 إلى 50 مليون سنة مضت.

\textbf{6.} $(20.00 - 5.59) \times 111.2/40 = 14.41 \times 111.2/40 \approx
\qty{40}{mm}$ في السنة.

\textbf{7.} $40/130 \approx 0.31$: أقل من الثلث.

\textbf{8.} من الخطوة 55--50 (146) إلى الخطوة 50--45 (89): نحو 50
مليون سنة مضت.

\textbf{9.} \omterm{def:g12:continental-crust:continental-crust}{القشرة القارية} أخف من أن تغوص: فحين بلغت
الخندق، لم يعد \omterm{def:g11:subduction-zones:subduction}{الاندساس} قادرًا على جرّ الهند بحرية، وأعاقتها القارة
الأخرى.

\textbf{10.} $(20.00 + 0.68) \times 111.2 = 20.68 \times 111.2 \approx
\qty{2300}{km}$.

\textbf{11.} $(29.65 - 21.15) \times 111.2 = 8.50 \times 111.2 \approx
\qty{945}{km}$.

\textbf{12.} $2300 - 945 \approx \qty{1354}{km}$.

\textbf{13.} اليوم $(29.65 - 20.00) \times 111.2 \approx \qty{1073}{km}$؛
وقبل 50 مليون سنة $(21.15 + 0.68) \times 111.2 = 21.83 \times 111.2
\approx \qty{2427}{km}$ (و$2427 - 1073 = 1354$).

\textbf{14.} $e = 1354/2427 \approx 0.56$، أي \qty{56}{\percent}.

\textbf{15.} \omterm{def:g12:mountain-building:fold}{الطيات}، \omterm{def:g12:mountain-building:thrust}{والصدوع المعكوسة}: \omterm{def:g12:mountain-building:thrust}{الصدوع الزاحفة} التي تحمل \omterm{def:g12:mountain-building:nappe}{الأغطية التكتونية}.

\textbf{16.} $H = 35 \times 2427/1073 \approx \qty{79}{km}$.

\textbf{17.} $\Delta H = 79 - 35 = \qty{44}{km}$؛ فسترتفع بمقدار
$44 \times 550/3380 \approx \qty{7.2}{km}$.

\textbf{18.} $\Delta H = 4.6 \times 3380/550 \approx \qty{28}{km}$، أي قشرة
سمكها $35 + 28 = \qty{63}{km}$، وهذا قريب من \qty{70}{km}
المقيسة.

\textbf{19.} القشرة لا تحتفظ بمقطعها العرضي: فأجزاء من آسيا
تُعصر نحو الشرق، والتعرية تزيل الصخر من أعالي
الجبال.

\textbf{20.} $2079 + 2300 \approx \qty{4380}{km}$: فمنذ نهاية
الطباشيري، قطعت الهند نحو \textbf{\qty{4400}{km} نحو الشمال}.
\end{solution}
